3. Econometric framework


The theoretical point of departure is the economic model of crime (Becker, 1968; Ehrlich, 1973), in which participation in illegal activity results from a rational comparison of the expected returns of legal and illegal opportunities, weighted by the probability and severity of sanctions. In this framework, a deterioration of the labor market lowers the opportunity cost of crime and should raise criminal activity (motivation effect), while falling income also reduces consumption opportunities for property crime (income effect); the net prediction is ambiguous and must be resolved empirically (Draca \& Machin, 2015). The empirical literature reviewed in Section 5 shows that the labor channel operates mainly on property crime, while homicide waves in Latin America respond to illicit markets, armed-group fragmentation, and state capacity. Against that backdrop, we estimate a deliberately modest set of specifications, each matched to what the available data can actually identify.



\subsection{Cross-sectional correlations (N = 24)}



Because provincial labor indicators exist for a single ENEMDU quarter (2024-II), the first block of evidence is descriptive: Pearson and Spearman correlations between the three provincial labor rates and provincial homicide rates, computed in the 2025 annual cross-section. With 24 provinces, these correlations are coarse descriptive statistics, not causal evidence; their role is to benchmark the raw association behind the labor market narrative (Chiricos, 1987, warns that cross-sectional correlations tend to overstate causal effects).



\subsection{Panel regressions: two-way fixed effects and the identification limit of a single labor cross-section}



The second specification is the canonical two-way fixed-effects (TWFE) model for province \textit{i} in quarter \textit{t}:



$$ h_{it} = \alpha_i + \lambda_t + \beta \, \ell_i + \varepsilon_{it}, $$



where $h_{it}$ is the homicide rate per 100,000, $\alpha_i$ and $\lambda_t$ are province and period fixed effects, and $\ell_i$ is the provincial labor indicator. The central identification problem is transparent: with a single labor cross-section, $\ell_i$ is time-invariant and collinear with the province fixed effects, so $\beta$ is not identified in the within dimension. The identification limit is structural, not incidental: the available data cannot separate the level association between labor conditions and homicides from permanent provincial heterogeneity. The fallback specification is therefore a cross-sectional OLS regression of (log) provincial homicide rates on labor indicators with regional dummies (coast, highlands, Amazon, and Gal\'apagos grouping) and province-clustered standard errors, exploiting cross-province variation within regions. Given only 24 clusters, clustered inference is complemented by the conservative t-distribution with G$-$1 degrees of freedom and, where feasible, wild cluster bootstrap p-values (Bertrand et al., 2004; MacKinnon \& Webb, 2017). These estimates are reported as associations; they are not causal parameters.



\subsection{Event study around Decreto 111 with a placebo in 2021Q1}



The third specification examines the timing of the state response. Decreto 111 (9 January 2024) declared an "internal armed conflict" and was followed by Decrees 218 (April 2024), 55 (July 2025), and 424 (June 2026). We estimate a dynamic event study (dynamic difference-in-differences) with province and period fixed effects, interacting relative-time indicators around 2024Q1:



$$ h_{it} = \alpha_i + \lambda_t + \sum_{k \neq -1} \delta_k \, D^k_{it} + \varepsilon_{it}, $$



where $D^k_{it}$ indicates that province $i$ is $k$ quarters away from the treatment quarter, with $k = -1$ normalized to zero. Because the decree was applied nationally at a single date, all provinces are treated simultaneously; in this design there are no "forbidden comparisons" between already-treated and not-yet-treated units, and the TWFE dynamic estimator is not subject to the heterogeneous-treatment-timing bias documented by de Chaisemartin \& D'Haultf{\oe}uille (2020) and Goodman-Bacon (2021); the interaction-weighted estimator of Sun \& Abraham (2021) coincides with the standard event study under a single treatment date (Callaway \& Sant'Anna, 2021; Borusyak et al., 2024). The event study is complemented by a placebo design assigning a pseudo-treatment to 2021Q1, before the escalation began. Two cautions govern its interpretation. First, the placebo is a falsification exercise of the estimation machinery, not a validation of parallel trends: pre-tests have low power, and conditioning the research design on their non-significance distorts inference (Roth, 2022). Second, SUTVA is an explicit risk given migration and commercial linkages among provinces, particularly between Guayas, Pichincha, and Esmeraldas.



\subsection{Synthetic control for Esmeraldas}



The fourth specification asks what would have happened to a high-lethality province in the absence of the post-2024 security regime. The synthetic control method (SCM) constructs the counterfactual of Esmeraldas as a convex combination of untreated donor provinces that replicates its pre-treatment trajectory in the outcome and in predictors (Abadie \& Gardeazabal, 2003; Abadie et al., 2010, 2015). The donor pool comprises the other 23 provinces; the pre-treatment window covers 2018Q1--2023Q4 and the post-treatment window 2024Q1--2026Q2. Inference relies on placebo permutation -- comparing the post-treatment gap of Esmeraldas with the distribution of gaps obtained by treating each donor province fictitiously -- which, with 24 units, yields a minimum achievable placebo p-value of approximately 1/24. The reported statistics (pre-treatment RMSE of 4.78 and mean post-treatment gap of $-$4.87 homicides per 100,000) are therefore illustrative of a single case, not fine-grained inference, and must not be read as an evaluation of any specific intervention in Esmeraldas (Abadie et al., 2010).



\subsection{Methods explicitly deferred: staggered DiD, Double ML, and related tools}



Several robust methods are part of the project's roadmap but are not applicable to the current data, and we document why rather than deploy them inappropriately. Staggered difference-in-differences estimators (de Chaisemartin \& D'Haultf{\oe}uille, 2020; Callaway \& Sant'Anna, 2021; Sun \& Abraham, 2021; Borusyak et al., 2024) add power only if the decrees of 2024--2026 were adopted at different dates across provinces; with a single national treatment date they collapse to the standard DiD and add nothing. Double/debiased machine learning (Chernozhukov et al., 2018) and causal forests (Wager \& Athey, 2018) require large samples and rich covariates and would need individual-level ENEMDU microdata; they are not estimable on an aggregate panel of 816 observations and 24 clusters. Regression discontinuity designs require a credible eligibility threshold, which does not exist in Ecuador's security legislation (Imbens \& Lemieux, 2008). Instrumental variables require an exogenous shifter of crime or employment; no credible instrument has been identified, and the crime$\rightarrow$employment direction remains unidentifiable with current data (Angrist \& Pischke, 2009). The binding constraint for all of these methods is data: a quarterly provincial ENEMDU panel since 2018, canton-level crime and labor data, and interoperable administrative records.